\scanpage{58}
\mspage{14}
% archive sheet number: 57
C) \uline{Découpage} \deleted{\unsure{de}} suivant une partie \inserted{ouv. et f.} $\Sigma'\subset\Sigma_{0}$ de la sous-variété fermée des cercles de découpage --- cela correspond à une partie $A'$ de $A_{0}$ (arêtes \deleted{\unsure{ordin.}}\inserted{liées}), mais cette fois ci il ne s'agit pas de ``gommer'' i.e. de contracter en un point, mais de ``découper'', ce qui \deleted{\unsure{sur}} au niveau du graphe revient à découper une arête par son milieu (i.e. à la remplacer par deux arêtes libres).

Le \uline{marquage suit}. \deleted{Ici encore, pour \unsure{une} \illegible{1}} Mais il y a ici à distinguer entre le cas d'une opération élémentaire \inserted{en $C$} --- celle où \deleted{il y a \unsure{lieu} \illegible{2}} \inserted{$C$ est effectivement marqué i.e. $C\cap R\ne\unsure{\varnothing|\phi}$} i.e. $C\subset\deleted{\unsure{\widetilde{\Sigma}}}\Sigma^{+}$ --- \deleted{\unsure{qui}} à peu de choses près (une \unsure{ambiguité|ambiguïté} dans un ens. fini, \deleted{\unsure{)}} en plus) peut se \unsure{réverser|renverser}, à l'aide du produit de cette opération --- et celle où $C$ n'est pas marqué, Mais il semble que dans la façon dont nous allons travailler
